\documentclass[11pt, a4paper]{article}

% ── Packages ──────────────────────────────────────────────────────────────────
\usepackage[a4paper, top=2.5cm, bottom=2.5cm, left=2.5cm, right=2.5cm, headheight=15pt]{geometry}
\usepackage{fontenc}
\usepackage{inputenc}
\usepackage{microtype}
\usepackage{parskip}
\usepackage{titlesec}
\usepackage{titling}
\usepackage{fancyhdr}
\usepackage{xcolor}
\usepackage{enumitem}
\usepackage{mdframed}
\usepackage{booktabs}
\usepackage{hyperref}

% ── Colours ───────────────────────────────────────────────────────────────────
\definecolor{teal}{HTML}{1A6B6B}
\definecolor{lightteal}{HTML}{E8F4F4}
\definecolor{darkgray}{HTML}{1C1C1C}
\definecolor{midgray}{HTML}{555555}
\definecolor{rulegray}{HTML}{CCCCCC}

% ── Hyperref ──────────────────────────────────────────────────────────────────
\hypersetup{
  colorlinks=true,
  linkcolor=teal,
  urlcolor=teal,
  pdftitle={Nepal Health Evidence Initiative — Founding Proposal},
  pdfauthor={Founding Team}
}

% ── Section Headings ──────────────────────────────────────────────────────────
\titleformat{\section}
  {\Large\bfseries\color{teal}}
  {\thesection.}{0.6em}{}
  [\vspace{2pt}\textcolor{teal}{\titlerule[0.6pt]}\vspace{4pt}]

\titleformat{\subsection}
  {\normalsize\bfseries\color{darkgray}}
  {\thesubsection}{0.6em}{}

\titlespacing*{\section}{0pt}{18pt}{8pt}
\titlespacing*{\subsection}{0pt}{12pt}{4pt}

% ── List Styles ───────────────────────────────────────────────────────────────
\setlist[itemize,1]{
  label=\textcolor{teal}{\textbullet},
  leftmargin=1.4em,
  itemsep=3pt,
  topsep=4pt
}
\setlist[itemize,2]{
  label=\textcolor{midgray}{--},
  leftmargin=2.2em,
  itemsep=2pt,
  topsep=2pt
}
\setlist[enumerate,1]{
  label=\textcolor{teal}{\arabic*.},
  leftmargin=1.6em,
  itemsep=4pt,
  topsep=4pt
}

% ── Header / Footer ───────────────────────────────────────────────────────────
\pagestyle{fancy}
\fancyhf{}
\renewcommand{\headrulewidth}{0.4pt}
\renewcommand{\headrule}{\hbox to\headwidth{\color{teal}\leaders\hrule height \headrulewidth\hfill}}
\renewcommand{\footrulewidth}{0.4pt}
\renewcommand{\footrule}{\hbox to\headwidth{\color{teal}\leaders\hrule height \footrulewidth\hfill}}
\fancyhead[L]{\small\textit{\textcolor{midgray}{Nepal Health Evidence Initiative \quad|\quad Founding Proposal}}}
\fancyfoot[L]{\small\textit{\textcolor{midgray}{Confidential --- For Team Discussion}}}
\fancyfoot[R]{\small\textcolor{midgray}{\thepage}}

% ── Callout box ───────────────────────────────────────────────────────────────
\newmdenv[
  backgroundcolor=lightteal,
  linecolor=teal,
  linewidth=1.5pt,
  leftline=true, rightline=false, topline=false, bottomline=false,
  innerleftmargin=12pt,
  innerrightmargin=10pt,
  innertopmargin=8pt,
  innerbottommargin=8pt,
  skipabove=10pt,
  skipbelow=10pt
]{callout}

% ── Phase box ─────────────────────────────────────────────────────────────────
\newmdenv[
  backgroundcolor=white,
  linecolor=rulegray,
  linewidth=0.8pt,
  roundcorner=3pt,
  innerleftmargin=10pt,
  innerrightmargin=10pt,
  innertopmargin=7pt,
  innerbottommargin=7pt,
  skipabove=6pt,
  skipbelow=6pt
]{phasebox}

% ══════════════════════════════════════════════════════════════════════════════
\begin{document}

% ── Title Page ────────────────────────────────────────────────────────────────
\thispagestyle{empty}
\vspace*{3cm}

\begin{center}
  {\fontsize{28}{34}\selectfont\bfseries\color{teal}
    Nepal Health Evidence Initiative}\\[0.6em]
  {\large\itshape\color{midgray}
    Building the Evidence Infrastructure for\\[0.2em]
    Health Policy in Nepal}\\[1.2em]
  {\normalsize\color{midgray}
    Founding Proposal \quad\textbullet\quad March 2026}
\end{center}

\vspace{1.2em}
\begin{center}
  \textcolor{teal}{\rule{10cm}{0.6pt}}
\end{center}
\vspace{0.6em}

\begin{center}
  {\small\bfseries\color{teal} DRAFT --- FOR TEAM DISCUSSION ONLY}
\end{center}

\vspace{1em}
\begin{center}
  \textcolor{teal}{\rule{10cm}{0.6pt}}
\end{center}

\vspace{3cm}

\begin{callout}
\textbf{\textcolor{teal}{Purpose of this document}}\\[4pt]
This proposal outlines the founding vision, structure, and immediate action plan for the Nepal Health Evidence Initiative. It is shared with the founding team to invite critique, additions, and shared ownership. It is not a final plan --- it is the start of a conversation.
\end{callout}

\newpage

% ══════════════════════════════════════════════════════════════════════════════
\section{The Moment We Are In}

Nepal's March 2026 elections have produced one of the most dramatic political shifts in the country's recent history. The Rastriya Swatantra Party's surge --- built on the momentum of the 2025 Gen Z uprising --- reflects a deep public demand for accountability, transparency, and competent governance.

For the first time in years, there is a genuine appetite in Nepal's political landscape for evidence-based approaches to policy. A new generation of leaders is entering government without the entrenched patronage networks of the old parties. This is a narrow window: historically, new governments are most open to outside expertise in their first 12 to 18 months.

\begin{callout}
\textbf{We intend to use this window.}
\end{callout}

% ══════════════════════════════════════════════════════════════════════════════
\section{Mission and Focus}

The \textbf{Nepal Health Evidence Initiative} (working name) is a health policy research network dedicated to one thing: making rigorous evidence accessible and actionable for Nepal's health policymakers, practitioners, and public.

We focus \textbf{exclusively on health}. This is a deliberate choice. By going deep rather than broad, we aim to become the definitive go-to resource for health policy evidence in Nepal --- not another generalist think tank that covers everything superficially.

Our core proposition is \textbf{translation}. Thousands of high-quality health economics studies are published every year. Almost none of them are synthesised, contextualised, and made accessible for Nepali policymakers. We fix that.

% ══════════════════════════════════════════════════════════════════════════════
\section{Founding Team}

\begin{itemize}
  \item \textbf{Sabin Subedi} --- PhD candidate in Economics at the University of Strathclyde, specialising in health economics and causal inference, with a dissertation focused on Nepal's National Health Insurance Programme. Sabin leads the technical research agenda and the AI-assisted database build.

  \item \textbf{Mukesh Adhikari} --- PhD student in Health Policy and Management at the University of North Carolina Gillings School of Global Public Health. Mukesh brings complementary expertise in health systems and policy, and co-leads the evidence synthesis and content strategy.

  \item \textbf{Prajjwal Dhungana} --- Field Coordinator, based in Kathmandu, currently working as a Research Assistant. Prajjwal is responsible for government liaison, consultation attendance, and on-the-ground institutional relationship management.

  \item \textbf{Suresh Tinkari} --- Advisor. IT Manager at the Health Insurance Board of Nepal. Suresh brings direct institutional knowledge of Nepal's health insurance system and provides a critical bridge between the initiative and the HIB.
\end{itemize}

% ══════════════════════════════════════════════════════════════════════════════
\section{What We Build}

\subsection{The Evidence Portal}

A publicly accessible, searchable database of global health policy research --- curated, structured, and contextualised for Nepal. Every paper in the database will be tagged across five dimensions:

\begin{itemize}
  \item Policy domain (e.g.\ health financing, maternal health, primary care)
  \item Intervention type
  \item Study design and evidence strength (RCT, quasi-experimental, descriptive)
  \item Country of study
  \item Nepal relevance score with a plain-English contextualisation note
\end{itemize}

Each featured paper will have a standardised \textbf{1-pager}: what was studied, what they found, and --- crucially --- what this would mean for Nepal specifically. This Nepal-specific translation layer is our primary comparative advantage.

\subsection{AI-Assisted Database Construction}

Building a high-quality evidence database has traditionally taken years of manual effort. We will use AI (Claude API) to parse research papers at scale and extract structured data: intervention, outcome, study design, country, key finding, evidence strength, and Nepal relevance. Sabin will take direct responsibility for building and maintaining this pipeline. Once the extraction is running, parsed entries will be shared with the full team for validity checks before anything is published.

The target is to have the initial database operational within one month of launch. This speed is only possible because of AI-assisted parsing --- it removes the manual bottleneck that has historically made evidence databases slow and resource-intensive to build.

\subsection{Automated Report Generation}

Once the database is live, we will build a report generation tool on top of it. Users --- government officials, INGOs, donors, researchers --- will be able to query a specific health topic and receive a structured synthesis report, drawing on papers from the database with evidence grading. All reports will be reviewed by a team member before release. The long-term model is a tiered service: free basic queries for the public, and paid commissioned reports for institutional clients such as FCDO, WHO Nepal, and UNICEF.

\subsection{The Blog}

A home for original insights and findings from our own research. Written in accessible language with a personal voice. \textbf{Named authors on every post} --- this is deliberate: contributors build a public profile and a CV line, which is how we attract and retain talented volunteers.

The blog will be published in \textbf{both Nepali and English}, making our work accessible to Nepali-speaking policymakers and practitioners as well as the international research community. This bilingual commitment is central to the initiative's identity.

Every post is peer-reviewed by another team member before publication. Given the effort required to write well in two languages and maintain rigorous quality, the blog is where the team will invest the most sustained effort. We publish when there is something genuinely worth saying, not on a forced schedule.

\subsection{Nepal Health Policy Digest}

A monthly newsletter distributed via Substack summarising the most relevant new global health research for Nepal. Five to eight papers per issue, each with a two-sentence summary and a one-line Nepal implication. Designed for busy policymakers who do not read journals.

This is our highest-leverage product for growing the mailing list quickly. It requires modest effort to produce and delivers consistent, recurring value to our audience.

% ══════════════════════════════════════════════════════════════════════════════
\section{Content Architecture}

The three content products serve distinct audiences and purposes:

\medskip
\begin{center}
\small
\begin{tabular}{p{2.8cm}p{4cm}p{4.5cm}p{2.2cm}}
\toprule
\textbf{Product} & \textbf{Audience} & \textbf{Purpose} & \textbf{Cadence} \\
\midrule
Blog & Researchers, students & Original findings, personal voice & As needed \\[4pt]
Digest & Policymakers, practitioners & Curated global evidence & Monthly \\[4pt]
Evidence Portal & Government, donors, INGOs & Structured synthesis, on-demand & Always-on \\
\bottomrule
\end{tabular}
\end{center}
\medskip

Substack is the subscription and distribution layer. All canonical content lives on our own domain. This ensures we own the reader relationship long-term.

% ══════════════════════════════════════════════════════════════════════════════
\section{Engagement Strategy}

\subsection{Government and Parliament}

\begin{itemize}
  \item Submit written evidence to Ministry of Health and Population consultations as they arise under the new government
  \item Connect with parliamentary research staff --- one sustained relationship here is worth more than a large social media following
  \item Target new government health policy contacts directly with a policy brief drawn from the NHIP dissertation findings
\end{itemize}

\subsection{Research Community}

\begin{itemize}
  \item Build and grow a mailing list from day one --- the Digest is the primary tool for this
  \item Partner with an established Nepali institution (IIDS, Policy Research Institute, or a university economics department) for local credibility and practical support
\end{itemize}

\subsection{International Networks}

\begin{itemize}
  \item Position the initiative as the Nepal node for global health economics networks (J-PAL, IGC, iHEA)
  \item Engage FCDO, WHO Nepal, UNICEF Nepal, and ADB as potential commissioned report clients once the evidence portal is live
\end{itemize}

% ══════════════════════════════════════════════════════════════════════════════
\section{Organisational Structure}

We will register a formal organisation in Nepal --- as an NGO or research institute --- to establish legal standing, procurement eligibility for government contracts, and donor funding access. Nepal registration is a prerequisite for working with most bilateral and multilateral funders (USAID, FCDO, GIZ).

We will also explore registering a parallel entity in the UK (or using an existing institutional home at Strathclyde) to maintain eligibility for UKRI and ESRC funding streams.

In the near term, operations will be lean: the two co-founders (Sabin and Mukesh), the field coordinator (Prajjwal), and the advisory support of Suresh. We do not build headcount ahead of funding.

% ══════════════════════════════════════════════════════════════════════════════
\section{Funding Pathway}

Funding strategy is something we want to discuss openly as a team, and everyone's ideas are welcome. The following are Sabin's initial thoughts to start that conversation.

\subsection*{Near-term possibilities}

\begin{itemize}
  \item \textbf{ESRC Impact Acceleration Account (IAA)} via University of Strathclyde --- this could be framed around research policy impact, with Nepal's new government as the enabling context. Sabin is planning to explore this with his supervisors.
  \item \textbf{British Academy International Partnerships Programme} --- a route already explored in a related context and worth revisiting for this initiative.
  \item In-kind support from a Nepali institutional partner for office space and local administrative costs.
\end{itemize}

\subsection*{Longer-term possibilities}

\begin{itemize}
  \item Commissioned synthesis reports for INGOs, bilateral donors, and UN agencies operating in Nepal
  \item Research grants from UKRI, Wellcome Trust, or global health foundations
  \item Consultancy arrangements with Nepal's Ministry of Health and Population
\end{itemize}

These are starting points. We will discuss the funding strategy in detail as a team and would welcome ideas from everyone.

% ══════════════════════════════════════════════════════════════════════════════
\section{Immediate Next Steps}

We propose the following actions within the next 30 days:

\begin{enumerate}
  \item Agree on the organisation name and register the domain
  \item Begin Nepal NGO registration process (Prajjwal to lead)
  \item Collectively announce the launch of the initiative --- website, social media, and mailing list simultaneously
  \item Initiate the ESRC IAA conversation with supervisors at Strathclyde
  \item Identify and approach one institutional partner in Nepal for a research collaboration MoU
\end{enumerate}

% ══════════════════════════════════════════════════════════════════════════════
\vspace{1.5em}
\textcolor{teal}{\rule{\linewidth}{0.6pt}}
\vspace{0.4em}

\textbf{\textcolor{teal}{A Final Note}}

\medskip
The window between a new government forming and being absorbed by institutional inertia is historically short --- 12 to 18 months. We have the credibility, the network, and the timing. What we need now is to move with intention.

This proposal is a starting point, not a final plan. We share it with the team to invite critique, additions, and ownership. The initiative will only work if everyone around this table believes in it and shapes it.

\vspace{0.4em}
\textcolor{teal}{\rule{\linewidth}{0.6pt}}

\end{document}